\subsection*{1.3. Completeness}
\hfill 13

\noindent\textbf{Theorem 1.8.} If $\mathcal{X}$ is a finite dimensional case, then all norms are equiv-
alent. That is, for given two norms $\|\cdot\|_1$ and $\|\cdot\|_2$ there are constants $m_1$
and $m_2$ such that
\[
\frac{1}{m_2} \|f\|_1 \leq \|f\|_2 \leq m_1\|f\|_1.
\tag{1.36}
\]

\noindent\textbf{Proof.} Clearly we can choose a basis $\{u_j\}_{j=1}^n$, and assume that $\|\cdot\|_2$ is
the usual Euclidean norm, $\left\|\sum_j a_j u_j\right\|_2 = \sum_j |a_j|^2$. Let $f = \sum_j a_j u_j$, then
by the triangle and Cauchy Schwartz inequalities
\[
\|f\|_1 \leq \sum_j |a_j| \|u_j\|_1 \leq \sqrt{\sum_j |a_j|^2} \sqrt{\sum_j \|u_j\|_1^2}
\tag{1.37}
\]
and we can choose $m_2 = \sqrt{\sum_j \|u_j\|_2^2}$.

In particular, if $f_n$ is convergent with respect to $\|\cdot\|_2$ it is also convergent
with respect to $\|\cdot\|_1$. Thus $\|\cdot\|_1$ is continuous with respect to $\|\cdot\|_2$ and attains
its minimum $m > 0$ on the unit sphere (which is compact by the Heine-Borel
theorem). Now choose $m_1 = 1/m$. $\square$

\noindent\textbf{Problem 1.5.} Show that $l^2(\mathbb{N})$ is a separable Hilbert space.

\noindent\textbf{Problem 1.6.} Let $s(f,g)$ be a skew linear form and $p(f) = s(f, \bar{f})$ the
associated quadratic form. Prove the parallelogram law
\[
p(f + g) + p(f – g) = 2p(f) + 2p(g)
\tag{1.38}
\]
and the polarization identity
\[
s(f,g) = \frac{1}{4} \left( p(f+g) - p(f - g) + ip(f – ig) – ip(f + ig) \right).
\tag{1.39}
\]

\noindent\textbf{Problem 1.7.} Show that the maximum norm does not satisfy the parallel-
ogram law.

\noindent\textbf{Problem 1.8.} Prove the claims made about $f_n$, defined in (1.35), in the
last example.

\subsection*{1.3. Completeness}

Since $L^2$ is not complete, how can we obtain a Hilbert space out of it? Well
the answer is simple: take the \textbf{completion}.

If $\mathcal{X}$ is a (incomplete) normed space, consider the set of all Cauchy
sequences $\mathcal{X}^\mathbb{N}$. Call two Cauchy sequences equivalent if their difference con-
verges to zero and denote by $\overline{\mathcal{X}}$ the set of all equivalence classes. It is easy
to see that $\mathcal{X}$ (and $\mathcal{X}^\mathbb{N}$) inherit the vector space structure from $\mathcal{X}$. Moreover,

\noindent\textbf{Lemma 1.9.} If $x_n$ is a Cauchy sequence, then $\|x_n\|$ converges.